\section{The full adelic image of the two elliptic frames}
\label{sec:adelic-cm-frames}

The two corrected elliptic frames detect no further relative abelian
characters on the profinite Grothendieck--Teichm\"uller group.  The
reason is an equality of full adelic images: the multiplier-one
subgroup and the cyclotomic Galois subgroup have the same image in the
product of the two CM tori.  We prove the integral comparison, identify
the one quadratic entanglement between the division fields, and obtain
a second normalization of the main character problem.

\subsection{Corrected frames and their integral targets}

Write $\mathcal G=\widehat{\mathrm{GT}}_0$ for the gentle group,
$G=\widehat{\mathrm{GT}}$, $H=G_{\mathbb Q}$, and
\[
 M_0=\ker(\lambda:\mathcal G\longrightarrow
                       \widehat{\mathbb Z}^{\times}),
 \qquad H^{\mathrm{cyc}}=G_{\mathbb Q^{\mathrm{ab}}}.
\]
We recall the normalization of the frames.  In the profinite braid
group put
\[
 \widehat B_3=\langle u,v\mid uvu=vuv\rangle,\qquad
 c=(uv)^3,\qquad s=uvu,\quad r=uv,
\]
\[
 \Gamma=\widehat B_3/\overline{\langle c\rangle}
             =\langle s,r\mid s^2=r^3=1\rangle.
\]
The free pure subgroup is
$F=\ker(\Gamma\to S_3)=\widehat F(x,y)$, where
$x=u^2$ and $y=v^2$.  The subgroup generated by these two elements
in $\widehat B_3$ identifies with $F$.  For $g\in\mathcal G$ its
braid action is
\[
 \Phi_g(u)=u^{\lambda_g},\qquad
 \Phi_g(v)=f_g^{-1}v^{\lambda_g}f_g,\qquad
 \Phi_g(c)=c^{\lambda_g},\qquad f_g\in F'.
\]
The unique pure frames satisfy
\[
 \Phi_g(s)=\beta_gs\beta_g^{-1},\qquad
 \Phi_h(r)=A_hrA_h^{-1}\quad(h\in M_0),
 \qquad \beta_g,A_h\in F.
\]
Their existence, uniqueness, and cocycles are as in
\cite[Lemmas~5 and~7]{H17}; in particular
\begin{equation}
 \beta_{gh}=\Phi_g(\beta_h)\beta_g,\qquad
 A_{gh}=\Phi_g(A_h)A_g\quad(g,h\in M_0),\qquad
 [\beta_g^{-1}]=(\rho_2(g),\rho_2(g)).
 \label{eq:adelic-frame-cocycles}
\end{equation}
The parity result gives $\rho_2(M_0)\subset2\widehat{\mathbb Z}$.

Set
\[
 U=\begin{pmatrix}1&1\\0&1\end{pmatrix},\quad
 L=\begin{pmatrix}1&0\\-1&1\end{pmatrix},\quad
 S=ULU=\begin{pmatrix}0&1\\-1&0\end{pmatrix},\quad
 R=-UL=\begin{pmatrix}0&-1\\1&-1\end{pmatrix}.
\]
Let $\widetilde q:\widehat B_3\to
\operatorname{SL}_2(\widehat{\mathbb Z})$ send $u,v$ to $U,L$,
and let $q$ be its restriction to $F$.  The intertwiner proved above
has the form
\begin{equation}
 D_g=\begin{pmatrix}\lambda_g&b_g\\0&1\end{pmatrix},\qquad
 \widetilde q\Phi_g=\operatorname{Ad}(D_g)\widetilde q,\qquad
 D_{gh}=D_gD_h.
 \label{eq:adelic-frame-intertwiner}
\end{equation}
Here all entries are adelic.  Its existence can also be seen directly,
prime by prime: if
$\widetilde q(\Phi_g(v))=\left(\begin{smallmatrix}
 a&b\\d&2-a\end{smallmatrix}\right)$, the braid equation gives
$d(1+\lambda_gd)=0$.  The zero possibility would make the two
braid images commute, contrary to their surjectivity onto
$\operatorname{SL}_2(\mathbb Z_\ell)$.  Thus
$d=-\lambda_g^{-1}$, $b=\lambda_g(a-1)^2$, and
$D_g=\left(\begin{smallmatrix}\lambda_g&\lambda_g(1-a)\\0&1
\end{smallmatrix}\right)$.  Its normalization makes it unique.

Put $O_3=\mathbb Z[\omega]$, where $\omega^2+\omega+1=0$,
and $O_4=\mathbb Z[i]$.  Embed these rings in matrices by
$\omega\mapsto R$ and $i\mapsto S$.  A bar denotes their quadratic
involution.  Define
\begin{equation}
 N_j=\{a\in(O_j\otimes\widehat{\mathbb Z})^{\times}:
                                      a\bar a=1\},\qquad
 N_j(4)=\{a\in N_j:a_2\equiv1\pmod4\}.
 \label{eq:adelic-cm-targets}
\end{equation}
The subscript $2$ in this formula means the component at the prime two.

\begin{lemma}\label{lem:adelic-corrected-frames}
The formulas
\begin{equation}
 T_g=q(A_g)^{-1}D_g,\qquad
 W_g=q(\beta_g)^{-1}D_g\qquad(g\in M_0)
 \label{eq:adelic-corrected-frames}
\end{equation}
define continuous homomorphisms
$T:M_0\to N_3(4)$ and $W:M_0\to N_4(4)$.
\end{lemma}

\begin{proof}
The centralizations are signed equalities.  Indeed the two lifts of
$\Phi_g(s)=\beta_gs\beta_g^{-1}$ in $\widehat B_3$ differ by
$c^e$.  Squaring, and using $s^2=c$ and $\Phi_g(c)=c$, gives
$2e=0$ in $\widehat{\mathbb Z}$, hence $e=0$.  Cubing the
corresponding identity for $r$, with $r^3=c$, similarly gives
$3e=0$.  Thus $T_g$ commutes with $R$ and $W_g$ commutes with $S$
as integral matrices, without a scalar ambiguity.

The centralizer rings of $R$ and $S$ in
$\operatorname{Mat}_2(\widehat{\mathbb Z})$ are respectively
$\widehat{\mathbb Z}[R]$ and $\widehat{\mathbb Z}[S]$.
For example $aI+bR$ has determinant $a^2-ab+b^2$, and
$aI+bS$ has determinant $a^2+b^2$.  Since $\det D_g=1$ on
$M_0$, this places $T_g,W_g$ in $N_3,N_4$.
Using \eqref{eq:adelic-frame-cocycles} and
\eqref{eq:adelic-frame-intertwiner} gives, with the order fixed,
\[
 q(A_{gh})^{-1}D_{gh}
 =q(A_g)^{-1}D_gq(A_h)^{-1}D_h=T_gT_h;
\]
the same calculation applies to $W$.

Modulo four the image of $F$ is the abelian four-group
\[
 V_4=\{I,U^2,L^2,U^2L^2\}.
\]
Consequently $q_4(f_g)=1$.  The braid covariance at $v$ makes
$D_g\bmod4$ commute with $L$, which forces the upper-right entry
of $D_g$ to be zero modulo four.  Since $q_4(A_g)\in V_4$
commutes with $R$ and $C_{V_4}(R)=\{I\}$, one has
$T_g\equiv I\pmod4$.  Finally the two abelian coordinates of
$\beta_g$ are $-\rho_2(g)$, both even.  The map $q_4$ factors
through their reductions modulo two, so $q_4(\beta_g)=I$ and
$W_g\equiv I\pmod4$ as well.  All maps used are continuous.
\end{proof}

\begin{lemma}\label{lem:adelic-kernel-normality}
The joint kernel of $T,W$ in $M_0$ is normal in $\mathcal G$.
\end{lemma}

\begin{proof}
The formula $q(\beta_g)^{-1}D_g$ defines a homomorphism to
$\operatorname{GL}_2(\widehat{\mathbb Z})$ for every $g\in\mathcal G$:
the calculation in the preceding proof uses the beta cocycle,
which is valid for every multiplier.  Thus its kernel is normal.

For the cubic frame, fix $g\in\mathcal G$, $h\in M_0$, and
$\eta=ghg^{-1}$.  Put $\epsilon\in\{1,-1\}$ equal to
$\lambda_g$ modulo three.  The torsion theorem gives a conjugator
$\gamma_g$ with
$\Phi_g(r)=\gamma_g r^\epsilon\gamma_g^{-1}$.
Because $\Phi_g$ induces the identity on $S_3$, this conjugator
may be chosen as $\gamma_g=B_gs^e$ with $B_g\in F$ and
$e\in\{0,1\}$.  Comparing
$\Phi_\eta\Phi_g(r)$ with $\Phi_g\Phi_h(r)$ and using the
uniqueness of the pure frame gives
\[
 A_\eta=\Phi_\eta(\gamma_g)^{-1}\Phi_g(A_h)\gamma_g.
\]
Lift $\gamma_g$ to $\widetilde\gamma_g=B_g(uvu)^e$ in
$\widehat B_3$.  Since $\lambda_\eta=1$, the equality
$\Phi_\eta(uvu)=\beta_\eta(uvu)\beta_\eta^{-1}$ is signed,
by the squaring argument above.  Also $uvu$ normalizes the embedded
free group.  Consequently the displayed expression, formed with
these braid lifts, belongs to that free group and is its unique
lift of $A_\eta$.  Writing
$Q_g=\widetilde q(\widetilde\gamma_g)$, covariance therefore yields
\[
 q(A_\eta)=D_\eta Q_g^{-1}D_\eta^{-1}
                 D_gq(A_h)D_g^{-1}Q_g.
\]
As $D_\eta=D_gD_hD_g^{-1}$, inversion and multiplication by
$D_\eta$ give
\begin{equation}
 T_\eta=(Q_g^{-1}D_g)\,T_h\,(Q_g^{-1}D_g)^{-1}.
 \label{eq:adelic-cubic-kernel-conjugation}
\end{equation}
Thus the kernel of $T$ in $M_0$ is also stable under every
conjugation by $\mathcal G$.  Intersecting the two kernels proves
the assertion.
\end{proof}

There are two parity maps on the target.  Let
$\pi_3=\sqrt{-3}$, an associate of $1-\omega$.  For $t\in N_3(4)$
and $w\in N_4(4)$ define
\begin{equation}
 (-1)^{\epsilon_3(t)}=t_3\bmod\pi_3\in\mathbb F_3^{\times},\qquad
 w_3\bmod3=S^a\quad\Longrightarrow\quad
 \epsilon_4(w)=a\bmod2.
 \label{eq:adelic-parities}
\end{equation}
The latter is well defined because the norm-one group in
$(O_4/3)^{\times}=\mathbb F_9^{\times}$ is cyclic of order four.
Both maps are surjective.  Write
\begin{equation}
 \mathcal E=\{(t,w)\in N_3(4)\times N_4(4):
                       \epsilon_3(t)=\epsilon_4(w)\}.
 \label{eq:adelic-equal-parity}
\end{equation}

\subsection{The CM division fields and the normalized images}

Let $K_j=O_j\otimes\mathbb Q$, and consider
\[
 E_3:Y^2=X^3+1,\qquad E_4:Y^2=X^3-X.
\]
Their adelic Tate representations are denoted $\rho_3,\rho_4$.
Over $K_j$ they take values in $(O_j\otimes\widehat{\mathbb Z})^
\times$.  We use the following precise CM input: the torsion
reciprocity theorem identifies the action modulo geometric
automorphisms with the ray class action
\cite[Theorem~12.27, printed p.~124]{MilneMF110}.
Since both rings have class number one, the ray class exact sequence
identifies the inverse limit with
$(O_j\otimes\widehat{\mathbb Z})^\times/O_j^\times$.
The action on the Tate module is abelian over $K_j$
\cite[Remark~12.12 and Proposition~12.13, printed p.~119]{MilneMF110}.

Put $\alpha=3^{1/4}$ and $r_0=(3/4)^{1/4}$, with positive real
roots, and define on $H^{\mathrm{cyc}}$
\[
 \kappa(\sigma)=\frac{\sigma(r_0)}{r_0}
               =\frac{\sigma(\alpha)}{\alpha}\in\{1,-1\}.
\]
The equality follows from $r_0=\alpha/\sqrt2$.

\begin{proposition}\label{prop:adelic-cm-images}
The full division fields are
\[
 \mathbb Q(E_j[\mathrm{tors}])=K_j^{\mathrm{ab}}.
\]
On $H^{\mathrm{cyc}}$, one has
\begin{equation}
 \rho_4(H^{\mathrm{cyc}})=N_4(4),\qquad
 \rho_3^{\circ}(H^{\mathrm{cyc}})=N_3(4),\qquad
 \rho_3^{\circ}=\kappa\rho_3.
 \label{eq:adelic-arithmetic-projections}
\end{equation}
The representations $\rho_3^{\circ}$ and $\rho_3$ have the same
kernel on $H^{\mathrm{cyc}}$.
\end{proposition}

\begin{proof}
First determine the images over the CM fields.  Put $\pi_4=1+i$.
The eight points of $E_4[\pi_4^3]$ are rational over $K_4$:
besides the four points of $E_4[2]$, they have coordinates
\[
 (i,\pm(1-i)),\qquad(-i,\pm(1+i)).
\]
For instance, with $[i](X,Y)=(-X,iY)$, the addition law gives
$(1+i)(i,1-i)=(1,0)$, then $(1+i)(1,0)=(0,0)$, and then zero.
The other three points give the same conclusion.  The eight points
exhaust the kernel because the endomorphism has degree eight.
The four roots of unity map bijectively to
$(O_4/\pi_4^3)^{\times}$.  Hence
\[
 C_4=\{u\in\widehat O_4^{\times}:u\equiv1\pmod{\pi_4^3}\}
       \longrightarrow\widehat O_4^{\times}/\mu_4
\]
is an isomorphism.  Rationality puts $\rho_4(G_{K_4})$ inside
$C_4$, and CM reciprocity makes its image modulo $\mu_4$ full.
Thus $\rho_4(G_{K_4})=C_4$.

For $E_3$, the two-torsion abscissae are
$-1,-\omega,-\omega^2$, and
$E_3[1-\omega]=\{O,(0,1),(0,-1)\}$.  These coprime kernels
make $E_3[2\pi_3]$ rational over $K_3$.
The reduction of $\mu_6$ is an isomorphism onto
\[
 (O_3/(2\pi_3))^{\times}
              \simeq\mathbb F_4^{\times}\times\mathbb F_3^{\times}.
\]
The same argument gives
\[
 \rho_3(G_{K_3})=C_3
       =\{u\in\widehat O_3^{\times}:u\equiv1\pmod{2\pi_3}\}.
\]

The CM theorem implies that the action on torsion points modulo
geometric automorphisms already cuts out $K_j^{\mathrm{ab}}$.
Since the full torsion action is
abelian over $K_j$, the full division field over $K_j$ is the same
field.  Moreover $K_3$ occurs among the two-torsion coordinates,
and $K_4$ occurs by the Weil pairing on four-torsion.  This proves
the displayed division-field identities over $\mathbb Q$.

Complex conjugation acts on $C_j$ by $u\mapsto\bar u$ and has
order two.  Thus the closed derived subgroup of the full Galois
image is $\{u/\bar u:u\in C_j\}$; this is already closed, as
the image of a compact group.  Kronecker--Weber identifies
$H^{\mathrm{cyc}}$ with the closed derived subgroup of $H$.
Consequently
\begin{equation}
 \rho_j(H^{\mathrm{cyc}})=\{u/\bar u:u\in C_j\}.
 \label{eq:adelic-ratio-image}
\end{equation}

At split primes the ratio map on units is
$(a,b)\mapsto(a/b,b/a)$ and is onto the local norm-one group.
At an unramified quadratic prime, Hilbert~90 supplies a preimage;
multiplication by a rational uniformizer power makes it a unit.
At the inert prime two for $K_3$, a norm-one unit $v\equiv1\pmod2$
has a unit preimage whose residue is fixed by the residue-field
involution.  Its residue therefore belongs to
$\mathbb F_2^\times=\{1\}$, so the preimage is also one modulo two.
At the ramified prime three, a Hilbert--90 preimage of a norm-one
unit $v\equiv1\pmod{\pi_3}$ has even $\pi_3$-valuation:
an odd valuation contributes the residue
$\pi_3/\bar\pi_3=-1$.  Multiplying by a rational power of three
and then by a sign makes it a unit congruent to one modulo $\pi_3$.
The converses follow by reduction.  Thus
\begin{equation}
 \rho_3(H^{\mathrm{cyc}})
 =\{v\in N_3:v_2\equiv1\pmod2,\quad
                     v_3\equiv1\pmod{\pi_3}\}=:M_3.
 \label{eq:adelic-unnormalized-cubic-image}
\end{equation}

For $K_4$ at two, a unit $u=a+bi\equiv1\pmod{\pi_4^3}$ has
$b$ even, so $u/\bar u\equiv1\pmod4$.  Conversely if
$v\bar v=1$ and $v\equiv1\pmod4$, then
$b_0=(1+v)/2$ is a unit congruent to one modulo two and
$b_0/\bar b_0=v$.  The two unit classes modulo $\pi_4^3$ that
are one modulo two are $1$ and $-1$.  Hence one of $b_0,-b_0$
belongs to $1+\pi_4^3O_{4,2}$.  Together with the other local
calculations, this proves $\rho_4(H^{\mathrm{cyc}})=N_4(4)$.

We now determine the cubic normalization without an unspecified
four-division comparison.  The norm-one units modulo four in
$O_3/4$ that are one modulo two are exactly $1,-1$:
write a unit as $1+2a$, and use
$\operatorname{Tr}_{\mathbb F_4/\mathbb F_2}(a)=0$.
Let $k=\mathbb Q^{\mathrm{ab}}$.  If
$Y^2=\prod_{i=1}^3(X-e_i)$ has all its roots in $k$, the halves
of $(e_i,0)$ are obtained from
\[
 a^2=e_i-e_j,\quad b^2=e_i-e_k,\qquad
 (X,Y)=(e_i+ab,\,ab(a+b)),
\]
with the possible signs.  Conversely the ratios $Y/(X-e_i)$
for the choices $(a,b)$ and $(-a,b)$ recover $a+b$ and $b-a$,
and hence $a,b$.  Thus its four-division field over $k$ is
generated by the square roots of the root differences.
For $E_3$ every nonzero root difference is a root of unity times
$\sqrt{-3}$.  Since $k$ contains every root of unity,
\begin{equation}
 k(E_3[4])=k(3^{1/4}).
 \label{eq:adelic-cubic-four-division}
\end{equation}
This extension is nontrivial: the real field
$\mathbb Q(3^{1/4})$ misses the nonreal conjugates of its
Eisenstein polynomial $X^4-3$, so it cannot lie in an abelian
extension of $\mathbb Q$.
It follows that the modulo-four sign of $\rho_3(\sigma)$ is
exactly $\kappa(\sigma)$.

Multiplication by this sign gives an isomorphism
$M_3\to N_3(4)$.  Its inverse sends $w\in N_3(4)$ to
$\varepsilon(w_3)w$, where
$\varepsilon(w_3)\in\{1,-1\}$ is its residue modulo $\pi_3$.
Indeed this makes the three-component one modulo $\pi_3$, while
the modulo-four sign at two becomes $\varepsilon(w_3)$.
The two maps are inverse homomorphisms.  This proves both the
normalized cubic image and the equality of kernels.
\end{proof}

\begin{lemma}\label{lem:adelic-cm-entanglement}
The intersection of the two full division fields is
\begin{equation}
 K_3^{\mathrm{ab}}\cap K_4^{\mathrm{ab}}
       =\mathbb Q^{\mathrm{ab}}(3^{1/4}).
 \label{eq:adelic-cm-intersection}
\end{equation}
In the representations of Proposition~\ref{prop:adelic-cm-images},
the joint image of $H^{\mathrm{cyc}}$ is exactly $\mathcal E$.
\end{lemma}

\begin{proof}
Put $L=K_3^{\mathrm{ab}}\cap K_4^{\mathrm{ab}}$.
It contains $\mathbb Q^{\mathrm{ab}}$ and is Galois over
$\mathbb Q$.  In $J=\operatorname{Gal}(L/\mathbb Q)$, the two
distinct index-two subgroups
$J_3=\operatorname{Gal}(L/K_3)$ and
$J_4=\operatorname{Gal}(L/K_4)$ are abelian.
Their intersection $Z$ is central, since $J_3,J_4$ generate $J$.
Choose $a\in J_3\setminus Z$ and $b\in J_4\setminus Z$.
Then $J=\langle Z,a,b\rangle$, $a^2,b^2\in Z$, and
$J'$ is generated by the central commutator $[a,b]$.
The identity $[a^2,b]=[a,b]^2=1$ makes $|J'|\leq2$.
Since $L$ contains the maximal abelian extension of $\mathbb Q$,
this proves $[L:\mathbb Q^{\mathrm{ab}}]\leq2$.

The splitting field $\mathbb Q(i,\alpha)$ of $X^4-3$ has dihedral
group of order eight.  Its subgroup over $\mathbb Q(i)$ is cyclic
of order four; its subgroup over $\mathbb Q(\sqrt{-3})$ is a
four-group.  It is therefore contained in both $K_j^{\mathrm{ab}}$.
The nontrivial extension in
\eqref{eq:adelic-cubic-four-division} gives the reverse bound and
proves \eqref{eq:adelic-cm-intersection}.

It remains to identify the common quadratic quotient in the stated
matrix coordinates.  By
\eqref{eq:adelic-unnormalized-cubic-image},
$\rho_3(\sigma)$ is one modulo $\pi_3$; hence
\[
 (-1)^{\epsilon_3(\rho_3^{\circ}(\sigma))}
                      =\kappa(\sigma).
\]
For $E_4$ its three-division polynomial is
$3X^4-6X^2-1$.  Choose a root $x$ with
$x^2=1+2\sqrt3/3$.  Then
\[
 \frac{x^2}{\sqrt3}=\frac{2+\sqrt3}{3}
     =\left(\frac{\sqrt6+\sqrt2}{2\sqrt3}\right)^2.
\]
Thus $3^{1/4}\in k(E_4[3])$.  The latter extension over $k$ is
cyclic of order four, since the reduction of $N_4(4)$ modulo
three is $\langle S\rangle$.  Its unique quadratic subextension
is consequently $k(3^{1/4})$, and
\[
 (-1)^{\epsilon_4(\rho_4(\sigma))}=\kappa(\sigma).
\]
The joint image of two Galois representations is the fibre product
of their separate images over the Galois group of the intersection
of their fields.  The separate images are full by
Proposition~\ref{prop:adelic-cm-images}; the intersection has order
two and the two quotient maps have just been computed.  The fibre
product is exactly \eqref{eq:adelic-equal-parity}.
\end{proof}

\subsection{A signed comparison on the Legendre family}

\begin{proposition}\label{prop:adelic-signed-path-comparison}
Choose the Tate-module bases by transport along paths to the
harmonic and equianharmonic fibres of the Legendre family.  For
$\sigma\in H^{\mathrm{cyc}}$ one then has the exact integral
identities
\begin{equation}
 W_\sigma=\rho_4(\sigma),\qquad
 T_\sigma=\rho_3^{\circ}(\sigma)
                       =\kappa(\sigma)\rho_3(\sigma).
 \label{eq:adelic-signed-path-comparison}
\end{equation}
In particular $(T,W)(H^{\mathrm{cyc}})=\mathcal E$.
\end{proposition}

\begin{proof}
Let $Y=\mathbb P^1\setminus\{0,1,\infty\}$ and consider the
Legendre family
\[
 \mathcal L_t:\quad Y_1^2=X_1(X_1-1)(X_1-t)
\]
with tangential base point at $t=0$ towards $1$.
Choose the integral vanishing-cycle basis in the peripheral
convention fixed by $x,y$.  The two local monodromies are
$U^2,L^2$: a loop about either collision of two ordered branch
points lifts to the square of the corresponding Dehn twist, and
the two vanishing cycles have intersection of absolute value one.
This identifies the geometric Tate-module representation with $q$
on $F$, at every prime.

At zero the special cubic is $Y_1^2=X_1^2(X_1-1)$, whose two
branches at the node have slopes $i,-i$.  Over
$\mathbb Q(i)((t))$ the reduction is split multiplicative.
Tate uniformization supplies a primitive cyclotomic line and a
free rank-one quotient in every integral Tate module; the only
additional character over $\mathbb Q((t))$ is the splitting
character of $\mathbb Q(i)$.
For an explicit uniformization over an arbitrary complete
nonarchimedean field, including the Galois compatibility and the
split multiplicative criterion, see
\cite[\S2, Theorem~2.2 and pp.~3--4]{PapikianTate}.
The integral filtration follows directly: in
$\overline K^{\times}/Q^{\mathbb Z}$, the $n$-torsion has kernel
$\mu_n$ and quotient $\mathbb Z/n$ under
$[z]\mapsto m\bmod n$ for $z^n=Q^m$.
The characteristic is zero, so this argument applies to every
prime, including two.

Both characters on the two graded pieces are trivial on
$H^{\mathrm{cyc}}$.  The tangential arithmetic matrix
$V_\sigma$ is therefore upper unipotent with both diagonal
entries one.  Functoriality gives
$q\Phi_\sigma=\operatorname{Ad}(V_\sigma)q$ on $F$.
Equation~\eqref{eq:adelic-frame-intertwiner} gives the same
covariance for $D_\sigma$.
An upper unipotent matrix commuting with $U^2$ and $L^2$ is the
identity: the second commutation gives twice its upper-right
entry equal to zero in each $\mathbb Z_\ell$.
Consequently
\begin{equation}
 V_\sigma=D_\sigma.
 \label{eq:adelic-tangential-matrix}
\end{equation}
This uniqueness uses the primitive integral filtration, so it
does not leave a sign undetermined at the prime two.

The geometric fundamental group of the quotient orbifold
$[Y/S_3]$ is $\Gamma$, with pure subgroup $F$.
The arithmetic action induced by the image of the chosen
tangential base point restricts to the canonical Ihara action
on $F$.  It is exactly $\Phi_\sigma$ on $\Gamma$.
Indeed two automorphisms of $\Gamma$ agreeing on $F$ differ,
on each group element, by an element of $C_\Gamma(F)$, and this
centralizer is trivial.  For an explicit argument, if $z$ centralizes
$F$, then $z^6\in F$ and $z^6$ lies in the centre of the nonabelian
free profinite group $F$, which is trivial.  Thus $z$ has finite
order.  A nontrivial torsion element of $\Gamma$ is conjugate into
one of its two finite factors, and its centralizer lies in that
factor \cite{HerfortRibes1985}.  Such a finite centralizer cannot
contain $F$, so $z=1$.

The transposition $t\mapsto1-t$ fixes $p_4=1/2$.
Choose an ordinary geometric path $\xi_4$ in $Y$ from the
tangential base point to $p_4$ so that its transported orbifold
isotropy is $s$.  Such a choice exists: every order-two subgroup
is conjugate into the order-two factor, and a conjugator preserving
the specified transposition in $S_3$ belongs to $F\langle s\rangle$.
It can therefore be replaced by a pure conjugator.
For the order-three isotropy use one of
$p_3=(1\pm\sqrt{-3})/2$ and choose an ordinary path $\xi_3$
whose transported positive generator is $r$.
The two points give the two possible pure conjugacy classes;
choose the one corresponding to $r$.  The same conjugator
argument, with centralizer $\langle r\rangle$ in $S_3$, then
makes the path adjustment pure.  Both chosen points and their
coordinate-permutation isotropies are fixed by $H^{\mathrm{cyc}}$.

Paths are composed as maps of fibre functors.  If $\xi$ is either
chosen path, set
\[
 c_\sigma=(\sigma\xi)^{-1}\xi\in F.
\]
For the isotropy $s_p$ at its endpoint, the identity
$s=\xi^{-1}s_p\xi$ gives
\[
 \Phi_\sigma(s)=(\sigma\xi)^{-1}s_p(\sigma\xi)
                        =c_\sigma s c_\sigma^{-1}.
\]
The analogous formula holds for $r$.  Uniqueness of the pure
frames yields $c_\sigma=\beta_\sigma$ or $A_\sigma$.

Now use the Tate sheaf on $Y$, where the Legendre family itself
is defined.  If $\rho_p(\sigma)$ is its matrix at the ordinary
point $p$, functoriality of parallel transport gives
\[
 q(\sigma\xi)=\rho_p(\sigma)q(\xi)D_\sigma^{-1}.
\]
In the transported basis put
$M_p(\sigma)=q(\xi)^{-1}\rho_p(\sigma)q(\xi)$.
Taking the inverse in the preceding formula and multiplying by
$q(\xi)$ proves
\begin{equation}
 q(c_\sigma)=D_\sigma M_p(\sigma)^{-1},\qquad
 q(c_\sigma)^{-1}D_\sigma=M_p(\sigma).
 \label{eq:adelic-path-order}
\end{equation}
The signed representation here is the one already fixed on $F$.
The calculation applies on every finite torsion sheaf, with
compatible transition maps.  It therefore proves an integral
adelic identity.  No descent of the signed Tate sheaf to
$[Y/S_3]$ is needed.

It remains to identify the two ordinary fibres.  At $p_4=1/2$,
the substitution $X_1=X+1/2$ gives
$Y_1^2=X^3-X/4$.  This is the quadratic twist of $E_4$ by
$1/2$, whose square root lies in $\mathbb Q^{\mathrm{ab}}$.
Its representation on $H^{\mathrm{cyc}}$ is therefore $\rho_4$.
At $p_3=(1+\sqrt{-3})/2$, substituting
$X_1=X+(1+p_3)/3$ gives
\[
 Y_1^2=X^3+\frac{\sqrt{-3}}9
             =X^3+\left(\frac1{\sqrt{-3}}\right)^3.
\]
This is the quadratic twist of $E_3$ by
$d=1/\sqrt{-3}$.  The quotient of a square root of $d$ by $r_0$
lies in $\mathbb Q^{\mathrm{ab}}$; its square is $-2i/3$, and
one may take $\sqrt{2/3}\,\zeta_8^{-1}$.
Its twisting character on $H^{\mathrm{cyc}}$ is consequently
$\kappa$.  The conjugate equianharmonic point changes $d$ by
$-1$ and gives the same character.

The transported CM algebras are the centralizers of $S$ and $R$
already computed in Lemma~\ref{lem:adelic-corrected-frames}.
Indeed Proposition~\ref{prop:adelic-cm-images} makes each local
arithmetic image open in its one-dimensional CM torus, so its
commutant over $\mathbb Q_\ell$ is the corresponding quadratic
algebra.  The matrices $S$ and $R$, which commute with that image,
generate its full quadratic integer ring.
At a prime, interchanging the two CM embeddings changes a unit
to its conjugate; it preserves the norm-one groups, their level-four
subgroups, and both parity maps.  Thus the CM image calculations
apply in these transported integral bases.
Equation~\eqref{eq:adelic-path-order} now gives
\eqref{eq:adelic-signed-path-comparison}, and
Lemma~\ref{lem:adelic-cm-entanglement} identifies the joint image.
\end{proof}

\subsection{Eliminating the entire adelic frame contribution}

Recall the global quadratic character $\Delta:\mathcal G\to
\mathbb F_2$ of Section~\ref{sec:global-frame-character}.
It is zero on $H$ and on $G$; on
$N_0=\ker(\lambda,\rho_2)$ it equals $a_3+a_2$.
The vanishing on $G$ uses the normalized frame theorem
\cite[Theorem~3]{H17}.  We require no stronger arithmetic input.

\begin{theorem}\label{thm:adelic-frame-elimination}
Put $G_1=G\cap M_0$.  Then
\begin{equation}
 (T,W)(G_1)=(T,W)(H^{\mathrm{cyc}})=\mathcal E.
 \label{eq:adelic-gt-arithmetic-image}
\end{equation}
In particular, the closed normal subgroup
\begin{equation}
 C_{\mathrm{CM}}=
 \ker\bigl((T,W):G_1\longrightarrow N_3(4)\times N_4(4)\bigr)
 \label{eq:adelic-cm-kernel}
\end{equation}
satisfies
\begin{equation}
 G=C_{\mathrm{CM}}H,\qquad
 C_{\mathrm{CM}}\cap H=
 G_{K_3^{\mathrm{ab}}K_4^{\mathrm{ab}}}.
 \label{eq:adelic-cm-normalization}
\end{equation}
The inclusion of $H$ induces the canonical isomorphism
\begin{equation}
 G/C_{\mathrm{CM}}\simeq
 \operatorname{Gal}(K_3^{\mathrm{ab}}K_4^{\mathrm{ab}}/\mathbb Q).
 \label{eq:adelic-cm-arithmetic-quotient}
\end{equation}
Every continuous abelian character of $G$ that is zero on $H$
and whose restriction to $G_1$ factors through $(T,W)$ is zero.
\end{theorem}

\begin{proof}
On $N_0$ the matrix $D_g$ is the identity modulo twelve.
The reductions of $q(A_g)$ and $q(\beta_g)$ therefore give
\[
 \epsilon_3(T_g)=a_3(g),\qquad
 \epsilon_4(W_g)=a_2(g).
\]
For the cubic group the first equality is reduction modulo
$\pi_3$: $R$ reduces to one and $-I$ to minus one.
For the harmonic group the second equality is exponent parity
in $\langle S\rangle$.  Inverting a frame does not change either
parity.  Thus
$(\epsilon_3\circ T)+ (\epsilon_4\circ W)=\Delta$ on $N_0$.

The arithmetic normalization by the Kummer cocycle of two gives
$M_0=N_0H^{\mathrm{cyc}}$: its value on $M_0$ is even, and its
image on $H^{\mathrm{cyc}}$ is exactly
$2\widehat{\mathbb Z}$.  On $H^{\mathrm{cyc}}$ both the sum of
the two parities and $\Delta$ are zero, by
Proposition~\ref{prop:adelic-signed-path-comparison} and the
arithmetic vanishing of $\Delta$.
All these maps are homomorphisms on $M_0$, so
\begin{equation}
 (\epsilon_3\circ T)+ (\epsilon_4\circ W)=\Delta
                          \quad\hbox{on }M_0.
 \label{eq:adelic-parity-delta}
\end{equation}
Since $\Delta|_G=0$, the image of $G_1$ lies in $\mathcal E$.
It contains the image of $H^{\mathrm{cyc}}$, which is already
all of $\mathcal E$.  This proves
\eqref{eq:adelic-gt-arithmetic-image}.

For $g\in G$, first choose $\sigma\in H$ with
$\lambda_\sigma=\lambda_g$.  Then $g\sigma^{-1}\in G_1$.
By \eqref{eq:adelic-gt-arithmetic-image}, choose
$h\in H^{\mathrm{cyc}}$ with the same corrected frame pair.
It follows that $g\sigma^{-1}h^{-1}\in C_{\mathrm{CM}}$, proving
$G=C_{\mathrm{CM}}H$.
The intersection formula follows from the exact signed comparison,
the equality of the normalized and unnormalized cubic kernels,
and the two full division-field identities.
Normality follows from Lemma~\ref{lem:adelic-kernel-normality}.
The compact surjection $H\to G/C_{\mathrm{CM}}$ has exactly the
kernel in \eqref{eq:adelic-cm-normalization}, proving
\eqref{eq:adelic-cm-arithmetic-quotient}.

Finally let $\chi$ be a character as in the statement.
Its restriction to $G_1$ is a character of the image
$\mathcal E$, and it is zero there because
$H^{\mathrm{cyc}}$ maps onto that image.  Thus $\chi|_{G_1}=0$.
Matching a multiplier by an element of $H$ then gives $\chi=0$
on $G$.
\end{proof}

In particular the same relative abelian quotient can be computed as
\[
 B\simeq C_{\mathrm{CM}}\big/
       \overline{[G,C_{\mathrm{CM}}]\,
                 G_{K_3^{\mathrm{ab}}K_4^{\mathrm{ab}}}}.
\]
For any prime $p$, its continuous $\mathbb F_p$-characters are the
$G$-invariant characters of $C_{\mathrm{CM}}$ that vanish on the
displayed Galois subgroup.  This follows from the general extension
criterion for a decomposition $G=NH$ with $N$ closed and normal.

This normalization concerns the full pair of CM frames.  Combining
it with simultaneous rational-radical normalization would require
the joint arithmetic image of the CM and radical representations;
the two separate surjectivity statements do not imply that joint
compatibility.
